% V. Експериментални резултати и анализ.
% Всяко число тук е измерено на машината от Табл. \ref{tab:machine}.
\section{Експериментални резултати и анализ}
\label{sec:results}

Коректността е предусловие за всичко останало и е проверена изцяло: 85 модулни случая без
сървър, 6 интеграционни срещу PostgreSQL 16.4 и 6 срещу MySQL 8.0.39, всички съвпадения, нула
разминавания. Модулните случаи включват контролните вектори на MD5, SHA-1, SHA-256,
HMAC-SHA-256, PBKDF2 и пълния обмен SCRAM-SHA-256 от RFC 7677~\cite{rfc7677}; интеграционните
проверяват, че записана стойност от тип \code{numeric}, \code{timestamp} и масив се прочита
обратно със същата стойност през двоичния формат.

\begin{table}[H]
\centering
\caption{Едно логическо действие за един ред, микросекунди, $n=500$ след 20 загряващи
повторения. Съобщенията и обръщанията са преброени от механизма за проследяване}
\label{tab:latency}
\begin{tabular}{L{3.8cm}L{1.6cm}L{1.5cm}L{1.4cm}L{1.4cm}L{1.6cm}L{1.5cm}}
\toprule
\textbf{Вид заявление} & \textbf{Медиана} & \textbf{95-и пр.} & \textbf{Средно} & \textbf{Ст. откл.} & \textbf{Съобщ.} & \textbf{Обръщ.} \\
\midrule
PostgreSQL, прост цикъл          & 612 & 700 & 625 & 39 & 5  & 1 \\
PostgreSQL, разширен, пропуск    & 636 & 733 & 651 & 45 & 12 & 1 \\
PostgreSQL, разширен, попадение  & 605 & 760 & 629 & 85 & 7  & 1 \\
MySQL, текстов протокол          & 643 & 725 & 650 & 41 & не е отчетено & не е отчетено \\
\bottomrule
\end{tabular}
\end{table}

Разсейването в този пуск е малко: стандартното отклонение е под петнадесет на сто от средната
стойност. Правилото се прилага буквално: където разликата е под едно стандартно отклонение,
извод не се прави. Първото такова място е самата Табл.~\ref{tab:latency}. Трите вида
заявление за PostgreSQL се различават с по-малко от петдесет микросекунди при стандартни
отклонения между 39 и 85, тоест измерването \emph{не} подкрепя твърдение кой от трите е
по-бърз. Причината е в средата: върху обратна връзка едно допълнително протоколно съобщение
струва почти нищо.

Десните две колони дават две наблюдения. Първо, и трите случая струват едно обръщане по
мрежата: разширеният цикъл \emph{не} плаща допълнително обръщане при първото изпълнение,
защото \code{Parse}, \code{Describe}, \code{Bind}, \code{Execute} и \code{Sync} се записват в
един буфер и се изпращат заедно. Хипотезата за две обръщания би била вярна само за клиент,
който изпраща \code{Flush} и чака отговор след \code{Parse}. Второ, кешът намалява
съобщенията от 12 на 7, тоест с 42 на сто: спестени са \code{Parse} и \code{Describe} в
изпратената посока и \code{ParseComplete}, \code{ParameterDescription} и
\code{RowDescription} в получената. Тези пет съобщения не се превръщат в измеримо време тук,
затова изводът е по-тесен от очаквания: дали спестените съобщения стават спестено време
зависи от цената на едно съобщение по мрежата, а обратната връзка е най-неблагоприятният
случай за кеша.

\begin{table}[H]
\centering
\caption{Резултат от 5000 реда и четири колони, микросекунди, $n=100$. Последната колона е
времето само в декодиращите функции, измерено с отделен часовник}
\label{tab:formats}
\begin{tabular}{L{4.2cm}L{1.9cm}L{1.7cm}L{1.7cm}L{1.7cm}L{2.4cm}}
\toprule
\textbf{Случай} & \textbf{Медиана} & \textbf{95-и пр.} & \textbf{Средно} & \textbf{Ст. откл.} & \textbf{Декодиране} \\
\midrule
PostgreSQL, текстов & 3638 & 4479 & 3817 & 585 & 1846 \\
PostgreSQL, двоичен & 3115 & 5117 & 3333 & 715 & 250 \\
MySQL, текстов      & 5585 & 5991 & 5647 & 304 & не е отделяно \\
\bottomrule
\end{tabular}
\end{table}

Разликата в декодирането е 7,4 пъти: 1846 срещу 250 микросекунди за 5000 реда по три
декодирани колони, тоест 123 срещу 17 наносекунди на стойност. Причината е в кода: двоичното
декодиране на цяло число е едно зареждане и разместване на байтове, докато текстовото минава
през разбор на десетичен запис, а двоичното декодиране на \code{double} е
\code{std::bit\_cast}, докато текстовото вика \code{strtod}. Общото време също е по-малко за
двоичния формат, но печели \emph{по-малко}, отколкото декодирането спестява: 523 микросекунди
срещу 1596. Клиентът декодира ред, докато сървърът още произвежда следващите, тоест част от
спестеното време освобождава време, което така или иначе щеше да се изчака.

\begin{table}[H]
\centering
\caption{Закъснителност на заемане на връзка от пул, микросекунди}
\label{tab:pool}
\begin{tabular}{L{1.8cm}L{2.2cm}L{2.0cm}L{2.0cm}L{2.0cm}L{2.2cm}}
\toprule
\textbf{Размер} & \textbf{Едновремен.} & \textbf{Медиана} & \textbf{95-и пр.} & \textbf{Ст. откл.} & \textbf{Чакали} \\
\midrule
1 & 1  & 2,1  & 3,5   & 2,9 & 0 от 40 \\
1 & 4  & 1957 & 2213  & 245 & 159 от 160 \\
1 & 16 & 9570 & 10363 & 956 & 639 от 640 \\
2 & 1  & 0,6  & 1,2   & 0,2 & 0 от 40 \\
2 & 4  & 749  & 904   & 121 & 158 от 160 \\
2 & 16 & 5210 & 6384  & 650 & 638 от 640 \\
4 & 1  & 0,5  & 0,7   & 0,2 & 0 от 40 \\
4 & 4  & 0,5  & 1,2   & 0,5 & 0 от 160 \\
4 & 16 & 2343 & 2630  & 297 & 636 от 640 \\
8 & 1  & 0,4  & 0,8   & 0,1 & 0 от 40 \\
8 & 4  & 0,4  & 0,9   & 0,5 & 0 от 160 \\
8 & 16 & 808  & 1028  & 165 & 630 от 640 \\
\bottomrule
\end{tabular}
\end{table}

Когато едновременността не надхвърля размера на пула, заемането струва под три микросекунди и
нито едно заемане не изчаква. Когато я надхвърли, заемането се превръща в чакане на опашка и
цената му е времето за изпълнение на чуждо заявление, тоест хиляди микросекунди вместо части
от една. Разликата между двата режима е три порядъка и е най-голямата, която измерването
откри. При едновременност 16 медианата пада от 9570 микросекунди при пул от една връзка до
808 при пул от осем, тоест осемкратното увеличение намалява чакането 11,8 пъти, повече от
пропорционално: при по-голям пул по-малка част от заявките изобщо стигат до опашката, което
личи и от последната колона. Стандартното отклонение при насищане е около една пета от
медианата или по-малко, тоест величината е чакане от самия пул, а не шум от средата.

Разликите се отнасят към три източника. \textbf{Брой копия}: двоичният формат печели 7,4 пъти
в декодирането, разликата е изцяло в клиента и е точно това, което прякото говорене по
протокол прави наблюдаемо, а именно един ред в Табл.~\ref{tab:formats} вместо скрито
поведение в чужда библиотека. \textbf{Брой съобщения}: кешът спестява 42 на сто, точно по
броене, но съответната разлика във времето е под разсейването и затова не се твърди.
\textbf{Режим на работа}: заемането от пул има два режима с три порядъка разлика, а границата
е едновременност, равна на размера на пула; това е единственият показател, при който изборът
на параметър променя цената качествено.

Четири неща не са измерени. Сравнение с \code{libpq} и \code{libmysqlclient} липсва изцяло,
защото те не са налични на машината; програмата го съобщава и не отчита число, а твърдение за
отношението към доставчиковите библиотеки не се прави. Върховата потребена памет не се
отчита, затова колона за памет липсва във всички таблици. Измерване при мрежа с ненулево
закъснение не е правено, а то е единственият начин да се провери дали спестените 42 на сто
съобщения се превръщат в спестено време. Измерена е закъснителността на заемане от пул, но не
и пропускливостта, тоест броят завършени заявления за единица време.

\TODO{сравнение с \code{libpq} и \code{libmysqlclient} на машина, на която те са инсталирани;
клонът в измервателната програма съществува и чака само средата}

\TODO{измерване на върховата потребена памет}
